% Source: https://artofproblemsolving.com/community/c3h1263295_plane_confusion
% Stable identity: aops-topic-1263295-q01
\begin{asy}
usepackage("amsmath");
size(4.5cm);
defaultpen(fontsize(10pt));
import three; currentprojection=orthographic(5,3,2);
draw(surface((-1,-1,0)--(1,-1,0)--(1,1,0)--(-1,1,0)--cycle),palegray+opacity(0.5));
draw((-1,-1,0)--(1,-1,0)--(1,1,0)--(-1,1,0)--cycle);
draw(surface((-1,0,0)--(1,0,0)--(1,0,1.4)--(-1,0,1.4)--cycle),lightblue+opacity(0.5));
draw(surface((0,-1,0)--(0,1,0)--(0,1,1.4)--(0,-1,1.4)--cycle),pink+opacity(0.5));
draw((0,0,0)--(0,0,1.6),red+linewidth(1.2));
label("$P$",(0.9,0.9,0)); label("$M$",(-0.8,0,1.2)); label("$N$",(0,0.8,1.2)); label("$m$",(0,0,1.6),N);
label("\mbox{}",(min(currentpicture,user=true)+max(currentpicture,user=true))/2,black); // keeps the text after the figure black when LaTeX typesets the labels inline
\end{asy}

Let $a=M\cap P$ and $b=N\cap P$. If $a$ and $b$ were parallel (or equal), then $M$ and $N$, both perpendicular to $P$ through parallel lines of $P$, would be parallel or equal, contradicting that they meet in a single line. So $a$ and $b$ meet at a point $O$ of $P$.

Let $\ell$ be the line through $O$ perpendicular to $P$. Since $M\perp P$ and $M$ contains the point $O$ of $P$, the perpendicular to $P$ at $O$ lies in $M$ (a plane perpendicular to $P$ contains the perpendicular to $P$ at each of its points on $P$). Likewise $\ell$ lies in $N$. Hence $\ell\subseteq M\cap N=m$, so $m=\ell$ is perpendicular to $P$. $\blacksquare$

(Intuitively: two walls standing straight up on a floor meet in a vertical edge.)
